\chapter{付録：LaTeX 記述チートシート}
\label{ch:latex_cheatsheet}

本付録は，チーム内で LaTeX に不慣れなメンバーが最低限書けるようになるための早見表である。必要に応じて削除・拡張して使うこと。

\section{章・節・小節}

\begin{lstlisting}[language=TeX]
\chapter{章タイトル}       % 章（jex では大きな区切り）
\section{節タイトル}       % 節
\subsection{小節タイトル}  % 小節
\subsubsection{小々節}     % さらに細かい
\end{lstlisting}

ラベル＆相互参照：

\begin{lstlisting}[language=TeX]
\section{同期方式}\label{sec:sync}
... 詳細は \ref{sec:sync} 節を参照 ...
\end{lstlisting}

\section{表の書き方}

短い表は \verb|tabular| + \verb|booktabs|，横幅を揃えたい場合は \verb|tabularx|，ページをまたぐ長い表は \verb|longtable| を使う。

\begin{lstlisting}[language=TeX]
\begin{table}[H]
  \centering
  \caption{サンプル表}\label{tbl:sample}
  \begin{tabular}{lcr}
    \toprule
    項目 & 数値 & 単位 \\
    \midrule
    遅延 & 5   & ms \\
    精度 & 99  & \% \\
    \bottomrule
  \end{tabular}
\end{table}
\end{lstlisting}

\section{図・画像の挿入}

\begin{lstlisting}[language=TeX]
\begin{figure}[H]
  \centering
  \includegraphics[width=0.7\textwidth]{fig/system_overview.pdf}
  \caption{システム全体構成}\label{fig:overview}
\end{figure}
\end{lstlisting}

本文から参照: 図\verb|\ref{fig:overview}|。
画像は \verb|fig/| ディレクトリに配置する。推奨形式は PDF / PNG。

\section{ソースコードの掲載}

\verb|listings| パッケージを使う。日本語コメントを含む場合は \verb|plistings| を併用すること（本雛形では読み込み済み）。

\begin{lstlisting}[language=C++,caption={3フェーズループ例},label={lst:3phase}]
void loop() {
  // 1. 入力フェーズ
  for (int i = 0; i < INPUT_COUNT; i++) {
    inputModules[i]->updateInput(data);
  }
  // 2. ロジックフェーズ
  applyPattern(data);
  // 3. 出力フェーズ
  for (int i = 0; i < OUTPUT_COUNT; i++) {
    outputModules[i]->updateOutput(data);
  }
}
\end{lstlisting}

\section{数式}

\verb|amsmath| を使う。インラインは \verb|$...$|，別行立ては \verb|\[...\]| または \verb|equation| 環境。

\begin{lstlisting}[language=TeX]
同期誤差の指標を式 \eqref{eq:jitter} に示す。
\begin{equation}
  \sigma_{\text{jitter}} = \sqrt{\frac{1}{N}\sum_{i=1}^{N}(t_i - \bar{t})^2}
  \label{eq:jitter}
\end{equation}
\end{lstlisting}

\section{箇条書き・脚注・引用}

\begin{lstlisting}[language=TeX]
\begin{itemize}
  \item 項目A
  \item 項目B\footnote{脚注はこう書く}
\end{itemize}

% 参考文献 (refs.bib) から引用
BLEの仕様\cite{ble_spec} に従い…
\end{lstlisting}

\section{よく使う記号・単位}

\begin{lstlisting}[language=TeX]
$\mu$s, $\approx$, $\pm$, $\sim$, $\times$, $\cdot$
% 単位付き: 10\,ms, 5\,V, 3.3\,V  (薄空白 \, を入れる)
\end{lstlisting}

\section{コンパイル手順}

\subsection{Dev Container（推奨）}
VSCode で本リポジトリを開き，\verb|Reopen in Container| を選ぶ。LaTeX Workshop が自動で \verb|latexmk| を走らせ，PDF を生成する。

\subsection{Docker 直接実行}
\begin{lstlisting}[language=bash]
cd report
docker run --rm \
  -v "$(pwd):/workspace" -w /workspace \
  ghcr.io/paperist/texlive-ja:debian \
  latexmk main.tex
\end{lstlisting}

\subsection{GitHub Actions}
\verb|report/**| 配下に変更のある PR / push で自動的に \verb|latexmk| が走り，生成された PDF が Artifact にアップロードされる。Actions タブから \verb|report-pdf| をダウンロードできる。
